\section*{第 \ref{ch:g2:born-grow-die} 章——出生、生长、死亡}

\begin{solution}{exo:g2:born-grow-die:1}
出生、幼体、\omterm{def:g2:born-grow-die:cycle}{成体}、\omterm{def:g2:born-grow-die:cycle}{老年期}（以死亡结束）。
\end{solution}

\begin{solution}{exo:g2:born-grow-die:2}
处于\omterm{def:g2:born-grow-die:cycle}{成体}\omterm{def:g2:born-grow-die:cycle}{阶段}。
\end{solution}

\begin{solution}{exo:g2:born-grow-die:3}
按从老到幼排列：祖母、母亲、小学生、
婴儿。
\end{solution}

\begin{solution}{exo:g2:born-grow-die:4}
樱桃种子发芽（出生）；幼树多年来不断长高；
成年树开花结樱桃——于是产生新的种子；
之后它变老并死亡。
\end{solution}

\begin{solution}{exo:g2:born-grow-die:5}
例如：婴儿长出牙齿，又换掉，长出更坚固的牙齿；
小鸡身上的黄色绒毛变成真正的\omterm{def:g1:animals-around-us:covering}{羽毛}；
幼树柔软的绿色茎变成坚硬的木质。
\end{solution}

\begin{solution}{exo:g2:born-grow-die:6}
按\omterm{prop:g2:born-grow-die:lifespan}{寿命}从短到长排列：蝴蝶（几周）、猫（十到
二十年）、橡树（数百年）。
\end{solution}

\begin{solution}{exo:g2:born-grow-die:7}
因为在道路走到尽头之前，\omterm{def:g2:born-grow-die:cycle}{成体}会产下幼体，
幼体又重新开始这段道路——这正是紫色箭头回指的地方。
一个生命是一段道路；世代相传的生命则像车轮一样转动。
\end{solution}

\begin{solution}{exo:g2:born-grow-die:8}
开放性答案。照片中通常显示婴儿、儿童、成人——
\omterm{def:g2:born-grow-die:cycle}{老年期}可能是尚未到来的\omterm{def:g2:born-grow-die:cycle}{阶段}。
\end{solution}

\begin{solution}{exo:g2:born-grow-die:9}
老树结的苹果里带有种子——
每颗被种下的种子都会长成一棵幼小的苹果树，即下一代。老树
死亡后，它的幼树延续这一循环，
花园里的苹果树便得以保留。
\end{solution}

\begin{solution}{exo:g2:born-grow-die:10}
不对——原因在于这颗鹅卵石\omterm{def:g1:living-or-not:nonliving}{从未有过生命}。只有\omterm{def:g1:living-or-not:living}{生物}
才会出生、生长和死亡；鹅卵石根本没有\omterm{def:g2:born-grow-die:cycle}{生命周期}，
因此“永不死亡”并不是力量，而是根本不活着。
\end{solution}
